\section{共享年表之二：大国均势与内部裂解}

\begin{event}{公元前597年}{邲之战}{可靠纪年}{邲}\eventanchor{SA-0597-BI}{春秋·邲之战}
\term{这一年发生了什么}：楚庄王击败晋军。《左传·宣公十二年》把晋军内部指挥分歧写成失败关键，也将楚庄王的克制塑造成战争伦理；经文则只留下更少的事件痕迹。
\term{后来改变了什么}：晋楚之间不存在永久霸主。中原政治从单一盟主转向大国反复均衡，诸侯和卿大夫有了更大的回旋空间。
\end{event}

\begin{event}{公元前546年}{弭兵之会}{可靠纪年}{宋}\eventanchor{SA-0546-MIBING}{春秋·弭兵之会}
\term{这一年发生了什么}：晋、楚等国在宋会盟，试图给长期战争设定边界。所谓“弭兵”不是现代意义的裁军：它更像承认双方势力范围、在资源耗竭后暂时冻结大规模冲突的外交安排。
\term{后来改变了什么}：会盟降低了部分直接冲突，却没有消除各国内部的权力争夺。外部均势越稳定，国内卿族越有机会积累资源。
\end{event}

\begin{event}{公元前506年}{柏举之战}{可靠纪年}{柏举}\eventanchor{SA-0506-BOJU}{春秋·柏举之战}
\term{这一年发生了什么}：吴军深入楚国，攻入郢都。《左传·定公四年》提供核心战役线；《史记》加入孙子等更强的个人叙事。
\term{后来改变了什么}：长江下游的吴进入中原政治舞台；楚虽受重创却未消失，说明一个国家的韧性并不只由一场战役决定。
\end{event}

\begin{debate}[孙子在柏举到底扮演了什么角色]
最早且最详细的《左传》柏举叙事没有提孙子；《史记》才把孙子放进吴国战略核心。银雀山汉简证明《孙子兵法》文本很早存在，却不能证明孙武必然直接指挥柏举。读者可以把孙子视作后世理解吴军胜利的重要传统，而不把它写成已经被春秋材料直接证实的事实。
\end{debate}

\begin{event}{公元前481年}{《春秋》记事终点}{传统分期}{鲁}\eventanchor{SA-0481-END}{春秋·记事终点}
\term{这一年发生了什么}：《春秋·哀公十四年》记“西狩获麟”，传统上以此年作为经文终点。
\term{后来改变了什么}：旧会盟秩序持续空洞化，国家开始更直接地重组户籍、军役和行政资源；战国不是突然开始，而是春秋后期问题的继续。
\end{event}

\begin{debate}[481年不是一条天然的历史分界线]
“获麟”在《公羊传》等传统中具有浓厚的象征意义；现代研究则会把经文停止传承、孔子及弟子传统和鲁国档案问题分开讨论。前481年是文本终点，战国的制度转变则是更长的过程，不能被一条年份线切断。
\end{debate}
